% sec3.tex -- writer W2. Section 3: the greedy chase and the end theorem.
\section{The greedy chase and the end theorem}\label{sec:end}

Throughout Sections~\ref{sec:end} and~\ref{sec:farnear}, $0<v\le\vz$, $\mu=v/(1-v)$, and $x$ is a working front
(Definition~\ref{def:flow}(b)) on $[0,U]$ from $\xz$ at $\ttd$, with $|\partial_u\xz|\equiv1$, so that a label is its arclength
coordinate at $\ttd$. For labels $a<b$ we write $\lambda:=b-a$ (the \emph{label length}) and call $(a,b)$ a \emph{pair}; for a pair
$J$ we write $\lambda_J,l(J),c(J),e_J,N(J)$.

\subsection{The chase}

\begin{definition}[Greedy chase]\label{def:chase}
Let $P$ be a $1$-Lipschitz path on $[0,\ttd]$ with $\xz(0)=P(\ttd)$, and let $b_{\mathrm c}\in(0,U]$ (the \emph{cut}). Put
$t_0:=\ttd$, $a_0:=0$. Given a leg start $t_j$ with $P(t_j)=x_{a_j}(t_j)$, let
\[
a_{j+1}:=\max\{u\in[a_j,b_{\mathrm c}]:\ l_{t_j}(a_j,u)\le\eps\}.
\]
If $a_{j+1}=b_{\mathrm c}$, the chase \emph{stops at} $t_j=:\Tstop$: leg $j+1$ is \emph{final}, it is not walked, and $P$ is
constant after $\Tstop$. Otherwise $P$ runs at unit speed along the ray from $P(t_j)$ through $x_{a_{j+1}}(t_j)$ until it meets
$x_{a_{j+1}}$, at time $t_{j+1}$; this requires $c(a_j,a_{j+1};t_j)>0$ (the \emph{proviso}). Leg $j+1$ is the interval
$[t_j,t_{j+1}]$, and $\uhat_{j+1}$ is its direction. The working front is computed with the path so extended.
\end{definition}

For $t\in(t_j,t_{j+1}]$, and for $t=t_0$ with $j=0$, we write $j(t):=j$; $a_{j(t)}$ is the \emph{base} at time $t$. A leg start
$t_j$, $j\ge1$, belongs to the legs $j$ and $j+1$: statements ``in leg $j+1$'' may be applied at $t_j$ with base $a_j$, and
statements ``in leg $j$'' at $t_j$ with base $a_{j-1}$; both are valid since $P(t_j)=x_{a_j}(t_j)$. Likewise, for a tail pair the
time $\tW$ is treated as a point of leg $\jW+1$, with base $a_{\jW}$ (Lemma~\ref{lem:local} and Section~\ref{sec:farnear}). Since $x(u,t)$ depends only on $\xz(u)$ and on $P|_{[\ttd,t]}$, the chase cut at
$b_{\mathrm c}$ depends only on $\xz|_{[0,b_{\mathrm c}]}$ and on the path up to the current time.

\begin{lemma}[The chase]\label{lem:chase}
Assume the proviso at every non-final leg start reached.
\begin{enumerate}
\item[(0)] Either the chase stops or $t_j\to\infty$; $P$ is $1$-Lipschitz.
\item[(1)] During leg $j+1$: $x_{a_{j+1}}=P+\rho\uhat_{j+1}$ with $\rho'=-1+vf(\rho)\le-(1-v)$, and $x_{a_j}=P-\sigma\uhat_{j+1}$ with
$\sigma\ge0$. So $x_{a_j},P,x_{a_{j+1}}$ are collinear in this order, and the leg lasts at most
$c(a_j,a_{j+1};t_j)/(1-v)\le\eps/(1-v)$.
\item[(2)] Every true label $u\in[a_j,a_{j+1}]$ has its direction seen by time $t_j$; this includes the final leg.
\item[(3)] During leg $j+1$: $l_t(a_j,a_{j+1})\le e^{2\mu}\eps$, $|P-x_{a_j}|\le e^{2\mu}\eps$, and $l_t(a,b)\le e^{2\mu}l_{t_j}(a,b)$
for all $a<b$.
\item[(4)] If a path $\tilde P$ equals $P$ on $[0,t_j]$, the working fronts of $P$ and $\tilde P$ agree on $[\ttd,t_j]$.
\end{enumerate}
\end{lemma}

\begin{proof}
By Lemma~\ref{lem:flow}(b) and Lemma~\ref{lem:field}(ii), $w_t=vw\,T\cdot AT\le(2v/\eps)w$ and, by Lemma~\ref{lem:excess}(i),
$l'\le(2v/\eps)l$.
(0) On $[\ttd,T]$, $w\le e^{2v(T-\ttd)/\eps}$. A non-final leg $j+1$ has $l_{t_j}(a_j,a_{j+1})=\eps$ (maximality and continuity of $l$
in the label), so if $t_j\le T$ it advances the base by at least $\eps e^{-2v(T-\ttd)/\eps}$; hence finitely many legs start
before $T$. $P$ is piecewise linear with speed $\le1$.
(1) Let $\rho$ solve $\rho'=-1+vf(\rho)$, $\rho(t_j)=c(a_j,a_{j+1};t_j)>0$. As $F(\rho\uhat)=f(\rho)\uhat$, $y:=P+\rho\uhat_{j+1}$ solves
$y'=vF(y-P)$ with $y(t_j)=x_{a_{j+1}}(t_j)$, so $y=x_{a_{j+1}}$ by uniqueness while $\rho>0$; $\rho'\le-(1-v)$ as $f\le1$. Likewise
$P-\sigma\uhat_{j+1}$ with $\sigma'=1+vf(\sigma)$, $\sigma(t_j)=0$ solves the equation of $x_{a_j}$ (note $F(-\sigma\uhat)=-f(\sigma)\uhat$),
and $\sigma\ge0$. Finally $c\le l_{t_j}(a_j,a_{j+1})\le\eps$.
(2) $\rho_u(t_j)=c(a_j,u;t_j)\le l_{t_j}(a_j,u)\le\eps$; apply Lemma~\ref{lem:seen}(b).
(3) The leg lasts at most $\eps/(1-v)$, so lengths grow at most by the factor $e^{(2v/\eps)\eps/(1-v)}=e^{2\mu}$; and
$|P-x_{a_j}|=\sigma\le c(a_j,a_{j+1};t)\le l_t(a_j,a_{j+1})$ by (1).
(4) Uniqueness for the equation of Definition~\ref{def:flow}.
\end{proof}

\subsection{Parameters and the end theorem}

\begin{definition}[Parameters]\label{def:param}
A \emph{parameter tuple} is
\[
p=(\gamma,R,m,s_0,\fF,\kappa,c_A,K_R,K_f,\Keps,\eta_1,\qc,\nu,\dL,\delta_1,A_0).
\]
Its \emph{derived constants} are
\begingroup\allowdisplaybreaks
\begin{align*}
&M:=mR^{-\gamma},\quad \chig:=\kappa\dL,\quad c_6:=\frac{1+\delta_1}{c_A},\quad \sigma_q:=\frac1{1-\qc^{-\gamma}},\quad
r_*:=\max\Bigl(\qc,\frac{2\qc}{\qc-1}\Bigr),\\
&\mathsf s(r):=\frac{\kappa^{1/2}}{c_A}\cdot\begin{cases}r/2,&1\le r\le2,\\(r-1)^{1/2},&r\ge2,\end{cases}\qquad
\omega_r:=\Bigl(1-\frac{m^2r}{4c_A}\Bigr)^{-1/2},\\
&\CT:=\omega_{r_*}\max\Bigl(\sigma_q\,\mathsf s(\qc),\,\mathsf s\Bigl(\frac{2\qc}{\qc-1}\Bigr)\Bigr),\quad \Cmu:=1+\kappa\CT^2,\\
&c_\rho^*:=c_A-\frac{2\delta_1}{K_f},\quad c_\eta:=c_A-\frac{\delta_1}{\Keps\eta_1},\quad c_R^*:=c_AK_R-\frac{\delta_1}{\Keps},\\
&K_1:=\frac{c_\rho^*K_f}{\kappa},\quad K_2:=\frac{c_R^*}{\kappa},\quad \camp:=\frac{\aD(K_1)\chig}{c_\rho^*},\\
&B_{2U}':=\frac{2\kappa^{3/2}(1+1/K_2)}{\pi},\quad
B_{2\eta}^*:=\frac{2\kappa^{3/2}(1+1/K_1)}{\pi c_\eta^2}\Bigl(\frac{\chig}{\eta_1}+\frac{\dL}{\Keps\eta_1^2}\Bigr),\\
&C_p':=\kappa^{1/2}\Bigl(1+\frac1{K_f}\Bigr)^{1/2+\gamma}K_f^{-1/2}+\kappa\CT\Bigl(1+\frac1{K_f}\Bigr)^\gamma,\quad
C_p^V:=(\kappa\Cmu)^{1/2}\Bigl(1+\frac1{K_f}\Bigr)^{1/2+\gamma},\\
&C_p^*:=\min(C_p',C_p^V),\quad C_{\mathrm{ar}}:=\frac{\kappa^{3/2}(1+1/K_f)^{3/2+\gamma}}{c_A},\quad
\hat C(x):=\min(C_p^*,C_{\mathrm{ar}}x^{1/2}),\\
&B_3':=\frac{2\kappa^{3/2}K_f^{-\nu}}{\pi(c_\rho^*)^3}\Bigl[2\dL K_f^{\nu+\gamma-2}\hat C(K_f)+\chig\int_{K_f}^\infty x^{\nu+\gamma-2}\hat C(x)\,dx\Bigr],\\
&\PsiI:=2+\log K_f,\quad\TI:=\dL(\kappa c_6+2)+c_6\chig\PsiI,\\
&B_5^*:=\frac1{c_\rho^*}\Bigl[\frac{C_p^*+\kappa^{1/2}K_f^{-1-\gamma}/\pi}{c_\rho^*}+K_f^{-\gamma}\Bigr]^2,\quad
I_5:=\frac{\chig(\eta_1+1/\Keps)^{2\gamma}}{2\gamma},
\end{align*}
where $\aD$ is the function of Lemma~\ref{lem:duhamel}; and $\bar\theta:=\fF^2+2\vz D'$, with the near-pair constant $D'$ of
Definition~\ref{def:Dprime} (it depends only on the tuple). With $x_*:=(C_p^*/C_{\mathrm{ar}})^2$ and $\hat p:=\nu+\gamma-2\in(-\frac32,-1)$,
the integral in $B_3'$ is $C_{\mathrm{ar}}(x_*^{\hat p+3/2}-K_f^{\hat p+3/2})/(\hat p+\frac32)+C_p^*x_*^{\hat p+1}/(-\hat p-1)$ if $x_*>K_f$
(then $\hat C(K_f)=C_{\mathrm{ar}}K_f^{1/2}$), and $C_p^*K_f^{\hat p+1}/(-\hat p-1)$ otherwise (then $\hat C(K_f)=C_p^*$). For a length $U_0>0$ and a
non-decreasing function $\Cfl:[0,\infty)\to[0,\infty)$ put $\rho_{\Cfl}(y):=\max(c_AK_RR,\Cfl(y))-\delta_1R/\Keps$ and, for $k=1,2$,
\begin{align*}
&J_k^{\Cfl}(U_0):=\frac{2\dL R}{\Keps(c_R^*R)^k}+\chig\Bigl[\int_{K_RR}^{\max(K_RR,U_0/2)}\rho_{\Cfl}^{-k}+\int_0^{U_0/2}\rho_{\Cfl}^{-k}\Bigr],\\
&J_k(U_0):=\min\Bigl\{\frac{\chig U_0}{(c_R^*R)^k},\,J_k^{\Cfl}(U_0)\Bigr\},\\
&\Xfar(U_0):=\vz\Bigl[J_1(U_0)+\frac1{c_\eta}\Bigl(\frac{\dL}{\Keps\eta_1}+\chig\log\frac{K_R}{\eta_1}\Bigr)\Bigr],\\
&\LJ(U_0):=\exp\bigl(\Xfar(U_0)+B_5^*I_5\vz m^2\bigr)+\vz\kappa\TI .
\end{align*}
\endgroup
\namedlabel{hyp:P}{\textup{(P0)--(P8)}}\textbf{Side conditions} for $(p,U_0,\Cfl)$:
\begin{itemize}
\item[(P0)] $1/2<\gamma$, $\gamma+\nu<1$; $0<R\le1$; $m>0$; $0<s_0\le\fF<1$; $1<\kappa\le4/3$; $0<c_A\le1$; $K_f\ge3$,
$K_f+1\le K_R$; $\Keps\ge2$; $0<\eta_1\le K_f/(K_f+1)$; $\qc>1$; $\nu\ge0$; $\dL,\delta_1,A_0\ge1$; $0<\vz<1$;
\item[(P1)] $1/(1-\vz)\le\dL$ and $e^{2\vz/(1-\vz)}\le\delta_1$;
\item[(P2)] $c_\rho^*K_f\ge2$, $c_R^*\Keps\ge1$, $c_6\le K_R\Keps$, $c_\eta>0$;
\item[(P3)] $\camp\vz\le\nu$;
\item[(P4)] $\vz\bigl[2\aD(K_1)\dL/(c_\rho^*K_f)+\aD(K_2)J_1(U_0)\bigr]\le\log A_0$;
\item[(P5)] $A_0\Bigl[s_0+\Bigl(\dfrac{K_R}{K_f}\Bigr)^{\!\nu}\dfrac{\vz}m\bigl(B_{2U}'RJ_2(U_0)+B_{2\eta}^*\bigr)+B_3'\vz\Bigr]\le\fF$;
\item[(P6)] $\bar\theta<1$;\qquad (P7) $\LJ(U_0)<\kappa$;\qquad (P8) $m^2r_*<4c_A$.
\end{itemize}
\end{definition}
By (P2), $c_\rho^*>0$ and $c_R^*>0$; so $\rho_{\Cfl}\ge c_R^*R>0$, and $J_k$ is non-decreasing in $U_0$. By (P1), $1/(1-v)\le\dL$
and $e^{2\mu}\le\delta_1$. Note $c_6\ge2$, $K_R\ge4$, that $\aD$ is decreasing, and that by (P8) $\omega_r$ is defined and
increasing for $1\le r\le r_*$, and $\mathsf s$ is increasing.

\medskip\noindent\namedlabel{hyp:end}{(H)}\textbf{Setting (H).} The working front on $[0,U]$ and the chase cut at $b_{\mathrm c}=U$;
a non-decreasing function $\Cfl:[0,\infty)\to[0,\infty)$;
$(p,U_0,\Cfl)$ satisfies (P0)--(P8); $U\le U_0$; $\eps\le R/\Keps$; an \emph{excuse length} $\We\ge0$ with $\We=0$ or
$U\ge2\We+R+\eps$; and a number $\delta_t\ge0$ (in the application, $\delta_t=e^{\vz K_w}-1$ of
Section~\ref{sec:dep}). (S1), (S1G) and (S2) are hypotheses at the time $\ttd$. A pair $(a,b)$, $0\le a<b\le U$, is a \emph{body pair} if $U-b\ge\We$, a \emph{tail pair} if $U-b<\We$, and
\emph{good} if $a\ge\We$ or $U-b\ge\We$. A pair that is not good has $\lambda>U-2\We\ge R+\eps$; so \emph{every pair with
$\lambda\le R$ is good}. The \emph{family of leg $j+1$} consists of the pairs with $\eps/2\le\lambda\le R$ and ($a=a_j$ or
$a\ge a_{j+1}$); the \emph{family} $\Fam(t)$ is the family of leg $j(t)+1$ (so $\Fam(t_j)$, $j\ge1$, is the family of leg $j$, which
contains that of leg $j+1$). The \emph{gap} of a pair of $\Fam(t)$ is $g:=a-a_{j(t)}\ge0$; the pair is \emph{far} if $g\ge K_f\lambda$
and \emph{near} otherwise. Let $\jW$ be the first $j$ with $a_j\ge\We$ and $\tW:=t_{\jW}$ (if the chase reaches it). The
\emph{admissible family} $\FamS(t)$ consists of the body pairs of $\Fam(t)$ if $t<\tW$ (or $\tW$ is not reached), and of all pairs
of $\Fam(t)$ if $t\ge\tW$. Hypotheses:
\begin{itemize}
\item[(S1)] $c(a,b;\ttd)\ge c_A\min(\lambda,K_RR)$ for all good pairs with $\lambda\ge\eps/2$;
\item[(S1G)] $c(a,b;\ttd)\ge\Cfl\bigl(\min(\lambda,a+U-b)\bigr)$ for all pairs with $\lambda>K_RR$;
\item[(S2)] $N(a,b;\ttd)\le s_0M\lambda^{1/2+\gamma}\min\bigl(1,K_f\lambda/\min(a,K_RR)\bigr)^\nu$ for all body pairs with
$\eps/2\le\lambda\le R$ (the last factor is $1$ if $a=0$);
\item[(TW)] (only if $\We>0$) if $\tW$ is reached, every tail pair with $\eps/2\le\lambda\le R$ satisfies, with $g_W:=a-a_{\jW}$,
\[
g_W\ge K_f\lambda,\quad l_{\tW}(a,b)\le(1+\delta_t)\lambda,\quad
N(a,b;\tW)\le s_0M\lambda^{1/2+\gamma}\min\Bigl(1,\frac{K_f\lambda}{\min(g_W,K_RR)}\Bigr)^{\nu};
\]
\item[(PT)] (only if $\We>0$) $(1+\delta_t)\LJ(U_0)<\kappa$.
\end{itemize}
If $\We=0$, then $\jW=0$, $\tW=\ttd$, there are no tail pairs, $\FamS=\Fam$, and (TW), (PT) are void. Families shrink in $t$ (a pair
of the family of leg $k+1$ has $a\ge a_k\ge a_{j+1}$ for $k>j$), and the gap of a fixed pair is non-increasing in $t$.

\begin{theorem}[End theorem]\label{thm:end}
In the setting \textup{(H)}, the proviso holds at every non-final leg start, and:
\begin{enumerate}
\item[(a)] \emph{(invariants)} for every $t\in[\ttd,\Tstop]$ and every pair of $\FamS(t)$,
\begin{equation}\label{eq:J}
\textup{(J1)}\quad l_t(a,b)\le\kappa\lambda,\qquad\qquad \textup{(J2)}\quad N(a,b;t)\le M\lambda^{1/2+\gamma};
\end{equation}
if $\We>0$, then $\tW$ is reached. Moreover, for all good pairs with $\lambda\ge\eps/2$ and all $t\ge\ttd$,
\begin{equation}\label{eq:CA}
\textup{(CA)}\qquad c(a,b;t)\ge c_A\min(\lambda,K_RR).
\end{equation}
\item[(b)] \emph{(the end)} every non-final leg advances the base by $a_{j+1}-a_j\in[\eps/\kappa,\eps]$; the chase stops at
$\Tstop\le\ttd+\kappa U/(1-v)$; the direction of every true label of $[0,U]$ is seen by $\Tstop$.
\item[(c)] \emph{(tilts)} at every $t\in[\ttd,\Tstop]$, for pairs $J\subset I$ of $\FamS(t)$: $|e_J-e_I|\le\CT m(\lambda_I/R)^\gamma$.
\item[(d)] \emph{(distance from the base)} for $t$ in leg $j+1$ and a label $s\ge a_j$ with ($U-s\ge\We$ or $a_j\ge\We$):
$\rho_s(t)\ge c_A\min(s-a_j,K_RR)-\delta_1\eps$ and $\max(\rho_s,\eps)\ge\min(s-a_j,K_RR)/c_6$; and, under \textup{(S1G)}, if
$s-a_j>K_RR$: $\rho_s(t)\ge\max\bigl(c_AK_RR,\Cfl(\min(s-a_j,a_j+U-s))\bigr)-\delta_1\eps$.
\item[(e)] \emph{(causality)} if a path agrees with the chase on $[0,t_j]$, then \textup{(a)}, \textup{(c)}, \textup{(d)} hold for it on
$[\ttd,t_j]$.
\end{enumerate}
No constant depends on $\eps$, $v$ or $U$.
\end{theorem}

The proof occupies the rest of this section and Section~\ref{sec:farnear}; it is completed in Section~\ref{sec:endproof}.

\subsection{Legs, gaps, distances, tilts}

\emph{Proof of \eqref{eq:CA}.} Chords are non-decreasing in $t$ (Lemma~\ref{lem:sep}(c)); apply (S1). In particular \eqref{eq:CA}
holds for every pair with $\eps/2\le\lambda\le R$ (such pairs are good, and $R\le K_RR$). \qed

\medskip
Let $T_{\mathrm{def}}\le\infty$ be the supremum of the times up to which the chase is defined (it is defined on $[\ttd,t_j]$ up to
each leg start $t_j$ reached).

\begin{lemma}[Legs]\label{lem:legs}
Let $t'\le T_{\mathrm{def}}$, and suppose \textup{(J1)} holds at every leg start $t_i\le t'$ for the pairs of $\FamS(t_i)$. Then for
every leg start $t_j\le t'$:
\begin{enumerate}
\item[(i)] if $U-a_j\ge\eps/2$, leg $j+1$ is final, or $a_{j+1}-a_j\in[\eps/\kappa,\eps]$ and $l_{t_j}(a_j,a_{j+1})=\eps$;
\item[(ii)] if $U-a_j<\eps/2$, leg $j+1$ is final; if $a_j<\We$, leg $j+1$ is not final;
\item[(iii)] if leg $j+1$ is not final, $c(a_j,a_{j+1};t_j)\ge c_A(a_{j+1}-a_j)>0$: the proviso holds;
\item[(iv)] at most $\kappa U/\eps$ non-final legs start in $[\ttd,t']$, each lasts at most $\dL\eps$, and every $t\le t'$ satisfies
$t-\ttd\le T_B:=\chig U$.
\end{enumerate}
The pairs used are $(a_j,s)$ with $a_j<s\le a_j+\eps$, and $(a_{j-1},U)$ in case \textup{(ii)}; they belong to $\FamS(t_j)$.
\end{lemma}

\begin{proof}
Always $a_{j+1}-a_j\le l_{t_j}(a_j,a_{j+1})\le\eps$, since $w\ge1$. \emph{Admissibility.} If $t_j<\tW$ then $a_j<\We$, and every
$s\le a_j+\eps$ has $U-s>U-\We-\eps\ge\We+R$: the pairs $(a_j,s)$ are body pairs; and $U-a_j>\eps$. If $t_j\ge\tW$,
$\FamS(t_j)=\Fam(t_j)$. Every pair $(a_j,s)$ with $\eps/2\le s-a_j\le R$ is in $\Fam(t_j)$ (for $j=0$ it has first point $a_0$; for
$j\ge1$, $a_j\ge a_j$ in the family of leg $j$).
(i) The label $s:=a_j+\eps/2$ is $\le U$, and $\eps/2\le R$ ($\Keps\ge2$); (J1) gives $l_{t_j}(a_j,s)\le\kappa\eps/2\le\eps$, so
$a_{j+1}\ge s$. If leg $j+1$ is non-final, $a_{j+1}<U$ and $l_{t_j}(a_j,a_{j+1})=\eps$ by maximality; the pair $(a_j,a_{j+1})$ has
$\eps/2\le\lambda\le\eps\le R$, and (J1) gives $\eps\le\kappa\lambda$.
(ii) A final leg $j+1$ has $U-a_j\le l_{t_j}(a_j,U)\le\eps$, so $a_j\ge U-\eps>\We$ if $\We>0$: legs with $a_j<\We$ are not final.
Now let $U-a_j<\eps/2$; then $t_j\ge\tW$. Induction on $j$. For $j=0$, $l_{t_0}(0,U)=U<\eps/2$. For $j\ge1$, leg $j$ was not final, so
$U-a_{j-1}\ge\eps/2$ by the induction hypothesis, and by (i) $\lambda_1:=a_j-a_{j-1}\in[\eps/\kappa,\eps]$ with
$l_{t_{j-1}}(a_{j-1},a_j)=\eps$. With $\lambda_2:=U-a_j<\eps/2$, the pair $(a_{j-1},U)$ has label length in $[\eps/2,3\eps/2]\subset
[\eps/2,R]$ and lies in $\Fam(t_j)$ (family of leg $j$). By (J1) and the monotonicity of lengths,
$l_{t_j}(a_j,U)\le\kappa(\lambda_1+\lambda_2)-\eps<\frac32\kappa\eps-\eps\le\eps$, as $\kappa\le4/3$; so $a_{j+1}=U$.
(iii) By (i), $\lambda:=a_{j+1}-a_j\in[\eps/2,R]$, and \eqref{eq:CA} gives $c\ge c_A\lambda$.
(iv) Each non-final leg starting in $[\ttd,t']$ advances the base by at least $\eps/\kappa$ (by (i), (ii)), and the advances add up to
at most $U$. Each leg lasts at most $\eps/(1-v)\le\dL\eps$ (Lemma~\ref{lem:chase}(1), (P1)). So $t-\ttd\le(\kappa U/\eps)\dL\eps=T_B$.
\end{proof}

\begin{lemma}[Gap--time integrals]\label{lem:gapint}
Let $t'$ be as in Lemma~\ref{lem:legs}, $a$ a label, and $\mathcal T\subset[\ttd,t']$ a set of times on which
$g(\tau):=a-a_{j(\tau)}\ge0$. The gap is constant on each leg and, by Lemma~\ref{lem:legs}(i), (ii), its values on distinct legs
differ by at least $\eps/\kappa$.
\begin{enumerate}
\item[(i)] If $h\ge0$ is non-increasing on $[g_{\min},\infty)$, where $g_{\min}$ is the smallest value $\ge g_0$ of $g$ on
$\mathcal T$, then
\[
\textstyle\int_{\{\tau\in\mathcal T:\,g(\tau)\ge g_0\}}h(g(\tau))\,d\tau\le\dL\eps\,h(g_{\min})+\chig\int_{g_{\min}}^\infty h(g)\,dg .
\]
\item[(ii)] If $h\ge0$ is non-decreasing on $[0,\bar g+\eps/\kappa]$ and $g\le\bar g$ on $\mathcal T$, then
\[
\textstyle\int_{\mathcal T}h(g(\tau))\,d\tau\le\chig\int_0^{\bar g+\eps/\kappa}h(g)\,dg .
\]
\item[(iii)] At most $\kappa x/\eps+1$ legs meet $\{\tau\in\mathcal T:g(\tau)<x\}$.
\item[(iv)] If $g(\tau)\in(g_0,Y]$ on $\mathcal T$, and $H\ge0$ is non-increasing on $(0,Y/2]$ and non-decreasing on $[Y/2,\infty)$, with
$H\le H_{\max}$ on $(g_0,Y]$, then \textup(an integral over an empty range is $0$\textup)
\[
\textstyle\int_{\mathcal T}H(g(\tau))\,d\tau\le2\dL\eps H_{\max}+\chig\bigl(\int_{g_0}^{Y/2}H(g)\,dg+\int_{Y/2}^YH(g)\,dg\bigr).
\]
\end{enumerate}
\end{lemma}

\begin{proof}
The bases of non-final legs increase by at least $\eps/\kappa$ (Lemma~\ref{lem:legs}), and each leg lasts at most $\dL\eps$.
(i) Let $g_1>\dots>g_n=g_{\min}$ be the values $\ge g_0$. For $i<n$, $g_i-\eps/\kappa\ge g_{i+1}$ and, $h$ being non-increasing,
$h(g_i)\le(\kappa/\eps)\int_{g_i-\eps/\kappa}^{g_i}h$; these intervals are disjoint and lie in $[g_{\min},\infty)$. So the integral is
$\le\dL\eps\sum_ih(g_i)\le\dL\eps\,h(g_{\min})+\dL\kappa\int_{g_{\min}}^\infty h$. (ii) Likewise $h(g_i)\le(\kappa/\eps)
\int_{g_i}^{g_i+\eps/\kappa}h$ for every value $g_i$, with disjoint intervals in $[0,\bar g+\eps/\kappa]$. (iii) The values in $[0,x)$
are spaced by at least $\eps/\kappa$. (iv) Let $g_1>\dots>g_n$ be the values of $g$ on $\mathcal T$. For each value $\le Y/2$ except
the smallest, $H(g_i)\le(\kappa/\eps)\int_{g_i-\eps/\kappa}^{g_i}H$ ($H$ is non-increasing on $(0,Y/2]$ and
$g_i-\eps/\kappa\ge g_{i+1}>g_0$); these intervals are disjoint and lie in $(g_0,Y/2]$. For each value $>Y/2$ except the largest,
$H(g_i)\le(\kappa/\eps)\int_{g_i}^{g_i+\eps/\kappa}H$ ($H$ is non-decreasing on $[Y/2,\infty)$ and $g_i+\eps/\kappa\le g_{i-1}\le Y$); these
intervals are disjoint and lie in $[Y/2,Y]$. The two exceptional values contribute at most $2\dL\eps H_{\max}$, the others at most
$\dL\eps(\kappa/\eps)(\int_{g_0}^{Y/2}H+\int_{Y/2}^YH)$.
\end{proof}

\begin{lemma}[Distance from the base]\label{lem:dist}
Let $t\le T_{\mathrm{def}}$ lie in leg $j+1$ \textup($t\in[t_j,t_{j+1}]$, or $t=t_j=\Tstop$\textup), and let $s\ge a_j$ be a label
with $U-s\ge\We$ or $a_j\ge\We$. Then
\begin{enumerate}
\item[(i)] $\rho_s(t)\ge c_A\min(s-a_j,K_RR)-\delta_1\eps$;
\item[(ii)] $\max(\rho_s(t),\eps)\ge\min(s-a_j,K_RR)/c_6$;
\item[(iii)] under \textup{(S1G)}, if $w:=s-a_j>K_RR$: $\rho_s(t)\ge\max\bigl(c_AK_RR,\Cfl(\min(w,a_j+U-s))\bigr)-\delta_1\eps$.
\end{enumerate}
We call (i) the first and (ii) the second form.
\end{lemma}

\begin{proof}
Let $w:=s-a_j$. If $w<\eps/2$, the first right side is negative ($\delta_1\ge1>c_A/2$). Otherwise $(a_j,s)$ is good, and by
\eqref{eq:CA} and $|P-x_{a_j}|\le e^{2\mu}\eps\le\delta_1\eps$ (Lemma~\ref{lem:chase}(3), (P1); at $t=t_j$, $P=x_{a_j}$):
$\rho_s\ge c(a_j,s;t)-|P-x_{a_j}|\ge c_A\min(w,K_RR)-\delta_1\eps$; here $c(a_j,s;t)\ge c(a_j,s;\ttd)$ (Lemma~\ref{lem:sep}(c)). Put
$y:=\min(w,K_RR)$. If $c_Ay-\delta_1\eps\ge y/c_6$ we are done; otherwise, as $c_A-1/c_6=c_A\delta_1/(1+\delta_1)$,
$y\,c_A\delta_1/(1+\delta_1)<\delta_1\eps$, i.e.\ $y<c_6\eps$. (iii) As in (i), $\rho_s\ge c(a_j,s;\ttd)-\delta_1\eps$, and
$c(a_j,s;\ttd)\ge\Cfl(\min(w,a_j+U-s))$ by (S1G) for the pair $(a_j,s)$.
\end{proof}

\begin{lemma}[Tilts]\label{lem:tilt}
Let $t\le T_{\mathrm{def}}$ lie in leg $j+1$, let $I=(a,b)$ be a pair of the family of leg $j+1$, and suppose that \textup{(J1)},
\textup{(J2)} hold at time $t$ for every pair of the family of leg $j+1$ contained in $[a_j,\max(b,a_{j+1})]$. Then:
\begin{enumerate}
\item[(o)] for pairs $J\subset I$ of the family of leg $j+1$ with $\lambda_I\le r\lambda_J$ and $1\le r<4c_A/m^2$:
$|e_J-e_I|\le\omega_r\mathsf s(r)M\lambda_I^\gamma$;
\item[(i)] for pairs $J\subset I$ of the family of leg $j+1$: $|e_J-e_I|\le\CT M\lambda_I^\gamma=\CT m(\lambda_I/R)^\gamma$;
\item[(ii)] for $a<s\le b$ and $\lambda_s:=\max(s-a,\eps/2)$,
\begin{equation}\label{eq:Emu}
\textup{(E-$\mu$)}\quad\int_a^s|T-e_{ab}|^2ds\le M^2\lambda^{2\gamma}\min(\Cmu\lambda_s,\lambda),\quad
\int_a^s|T-e_{ab}|\,ds\le M\lambda^\gamma\bigl(\kappa\lambda_s\min(\Cmu\lambda_s,\lambda)\bigr)^{1/2}.
\end{equation}
\end{enumerate}
\end{lemma}

\begin{proof}
All pairs below lie in $[a_j,\max(b,a_{j+1})]$ and in the family of leg $j+1$, so (J1), (J2) and \eqref{eq:CA} apply to them.
(o) Let $J\ne I$ (otherwise there is nothing to prove). By Lemma~\ref{lem:split}(i), $Z^2\le N(I)^2f(l(J))$ with
$f(\ell):=\ell(1-\ell/l(I))$; $f$ increases in $l(I)$, and $l(I)\le\kappa\lambda_I$, $l(J)\le\kappa\lambda_J$ by (J1). So
$f(l(J))\le\max_{0<\ell\le\kappa\lambda_J}\ell(1-\ell/(\kappa\lambda_I))=\kappa\lambda_J\tilde\rho(\rho)$, where $\rho:=\lambda_I/\lambda_J$
and $\tilde\rho(\rho)=\rho/4$ for $\rho\le2$, $\tilde\rho(\rho)=1-1/\rho$ for $\rho\ge2$. By (J2) and \eqref{eq:CA} ($c(J)\ge c_A\lambda_J$),
$Z/c(J)\le M\lambda_I^{1/2+\gamma}(\kappa\lambda_J\tilde\rho)^{1/2}/(c_A\lambda_J)=(\kappa^{1/2}/c_A)(\rho\tilde\rho(\rho))^{1/2}M\lambda_I^\gamma
=\mathsf s(\rho)M\lambda_I^\gamma\le\mathsf s(r)M\lambda_I^\gamma$, since $\rho\tilde\rho(\rho)$ is $\rho^2/4$ resp.\ $\rho-1$, $\mathsf s$ is increasing and
$\rho\le r$. Further $N(I)^2/(4c(J))\le M^2\lambda_I^{1+2\gamma}/(4c_A\lambda_J)\le m^2r/(4c_A)<1$ ($M\lambda_I^\gamma\le MR^\gamma=m$),
and Lemma~\ref{lem:split}(ii) gives the claim.

(i) Let $J=(a',b')\subset I$, $q:=\qc$, $\lambda_k:=\lambda_Iq^{-k}$ and $K:=\min\{k\ge1:\lambda_k\le\lambda_J\}$.
\emph{Case A: $a\ge a_{j+1}$, or $a'=a$, or $\lambda_I(1-1/q)\ge\eps$.} For $0\le k\le K-1$ put $\alpha_k:=\min(a',b-\lambda_k)$,
$I_k:=(\alpha_k,\alpha_k+\lambda_k)$, and $I_K:=J$. (1) $J\subset I_k\subset I$: $\alpha_k\ge\min(a',b-\lambda_I)=a$ and $\alpha_k\le a'$; the
right end is $\min(a'+\lambda_k,b)\le b$, and it is $\ge b'$ as $\lambda_k>\lambda_J$ for $k\le K-1$. (2) $\alpha_k$ is non-decreasing in $k$
and the right end non-increasing, so $I_{k+1}\subset I_k$. (3) Each $I_k$ is in the family: its length $\lambda_k\in(\lambda_J,\lambda_I]
\subset[\eps/2,R]$, and its first point is $a'$ (the first point of the family pair $J$) or $b-\lambda_k<a'$; in the second case $k\ge1$
($I_0=I$) and $b-\lambda_k\ge a+\lambda_I(1-1/q)$, which is $\ge a\ge a_{j+1}$ if $a\ge a_{j+1}$, and $\ge a_j+\eps\ge a_{j+1}$ if
$\lambda_I(1-1/q)\ge\eps$ ($a_{j+1}-a_j\le l_{t_j}(a_j,a_{j+1})\le\eps$, as $w\ge1$); if $a'=a$, then $\alpha_k=a$ for all $k$.
(4) $\lambda(I_k)/\lambda(I_{k+1})=q$ for $k+1\le K-1$, and $\lambda(I_{K-1})/\lambda_J\in[1,q]$. By (o) with $r=q$,
$|e_{I_{k+1}}-e_{I_k}|\le\omega_q\mathsf s(q)M\lambda_k^\gamma$ ($q\le r_*<4c_A/m^2$ by (P8)), and summing,
$|e_J-e_I|\le\omega_q\mathsf s(q)M\lambda_I^\gamma\sum_{k=0}^{K-1}q^{-k\gamma}<\omega_{r_*}\sigma_q\mathsf s(q)M\lambda_I^\gamma\le\CT M\lambda_I^\gamma$.
\emph{Case B: $a=a_j$, $a'>a$ and $\lambda_I(1-1/q)<\eps$.} Then $\lambda_I/\lambda_J<(\eps q/(q-1))/(\eps/2)=2q/(q-1)\le r_*$, and
(o) with $r=2q/(q-1)$ gives $|e_J-e_I|\le\omega_{r_*}\mathsf s(2q/(q-1))M\lambda_I^\gamma\le\CT M\lambda_I^\gamma$. The two cases are exhaustive.

(ii) $J:=(a,a+\lambda_s)$ is a pair of the family (first point $a$, length in $[\eps/2,\lambda]$) contained in $I$. By
Lemma~\ref{lem:static}(b) with $\uhat=e:=e_{ab}$, (J2), (J1) and (i),
$\int_a^s|T-e|^2\le\int_J|T-e|^2=N(J)^2+c(J)|e_J-e|^2\le M^2\lambda_s^{1+2\gamma}+\kappa\lambda_s\CT^2M^2\lambda^{2\gamma}\le\Cmu M^2\lambda^{2\gamma}\lambda_s$;
also $\int_a^s|T-e|^2\le\int_I|T-e|^2=N(I)^2\le M^2\lambda^{1+2\gamma}$. The second bound follows by Cauchy--Schwarz, as
$l(a,s)\le l(J)\le\kappa\lambda_s$.
\end{proof}

\begin{lemma}[Locality]\label{lem:local}
Let $t\le T_{\mathrm{def}}$, $(a,b)\in\FamS(t)$, and $\ttd\le t'\le t$, with $t'>\tW$ if $(a,b)$ is a tail pair. Then
\begin{enumerate}
\item[(i)] every pair of $\Fam(t')$ contained in $[a_{j(t')},\max(b,a_{j(t')+1})]$ belongs to $\FamS(t')$;
\item[(ii)] every $s\in[a,b]$ satisfies the hypothesis of Lemma~\ref{lem:dist} at $t'$ with the base $a_{j(t')}$; in particular the
pair $(a_{j(t')},s)$ is good if $s-a_{j(t')}\ge\eps/2$;
\item[(iii)] the pairs used by Lemma~\ref{lem:legs} at a leg start $t_i$ belong to $\FamS(t_i)$;
\item[(iv)] if $\We>0$: a tail pair with $\lambda\le R$ has $a>a_{\jW}$; the pairs of $\Fam(\tW)$ with first point $a_{\jW-1}$ are body
pairs.
\end{enumerate}
For a tail pair at $t'=\tW$, \textup{(i)} and \textup{(ii)} hold with $j(t')$ replaced by $\jW$ \textup($\tW$ belongs to leg
$\jW+1$, and $\FamS(\tW)=\Fam(\tW)$ contains the family of leg $\jW+1$\textup).
\end{lemma}

\begin{proof}
Note $a\ge a_{j(t)}\ge a_{j(t')}$. (i) If $t'\ge\tW$, $\FamS(t')=\Fam(t')$. If $t'<\tW$, then $(a,b)$ is a body pair and
$a_{j(t')}<\We$, so $a_{j(t')+1}<\We+\eps$; a pair $(a'',b'')$ in the interval has $U-b''\ge\min(U-b,U-\We-\eps)\ge\We$. (ii) For a
body pair $U-s\ge U-b\ge\We$; for a tail pair $t'>\tW$ gives $a_{j(t')}\ge a_{\jW}\ge\We$. (iii) is shown in the proof of
Lemma~\ref{lem:legs}. (iv) A tail pair with $\lambda\le R$ has $a\ge b-R>U-\We-R\ge\We+\eps>a_{\jW}$, since $a_{\jW}\le a_{\jW-1}+\eps$.
A pair with first point $a_{\jW-1}<\We$ and $\lambda\le R$ has $U-b>U-\We-R\ge\We$. The last sentence: $a_{\jW}\ge\We$.
\end{proof}

\noindent\textbf{Locality principle.} The proofs of Propositions~\ref{prop:far}--\ref{prop:length} for a pair $(a,b)\in\FamS(t)$
use the family bounds (J1), (J2) only for the pairs of Lemma~\ref{lem:local}(i), (iii) (at the times $t'$ allowed there),
\eqref{eq:CA} only for these pairs and for the pairs $(a_{j(t')},s)$, $s\in[a,b]$, of Lemma~\ref{lem:local}(ii) (through
Lemma~\ref{lem:dist}), the hypothesis (S1G) at $\ttd$ only for these pairs $(a_{j(t')},s)$ (through Lemma~\ref{lem:dist}(iii)),
and, beyond the bounds of the pairs of Lemma~\ref{lem:local}(i) (which include $(a,b)$ itself), only the start data of $(a,b)$ at
$t_\circ$. Each invocation in Section~\ref{sec:farnear} is of this kind.

\begin{lemma}[Continuity scheme]\label{lem:cont}
Consider the statement, for a time $t$:
\begin{quote}
\textup{(Imp)} every pair of $\FamS(t)$ satisfies $l_t(a,b)\le(1+\delta)\LJ(U_0)\lambda$ and
$E(a,b;t)\le\frac12\bar\theta M^2\lambda^{1+2\gamma}$, where $\delta:=0$ for body pairs and $\delta:=\delta_t$ for tail pairs.
\end{quote}
Suppose that, whenever $t^*\le T_{\mathrm{def}}$ is such that the chase has not stopped before $t^*$ and \textup{(J1)}, \textup{(J2)}
hold for all $t\in[\ttd,t^*]$ and all pairs of $\FamS(t)$, \textup{(Imp)} holds at every $t\in[\ttd,t^*]$. Then the proviso holds
at every non-final leg start, the chase stops, \textup{(J1)}, \textup{(J2)} hold for all $t\in[\ttd,\Tstop]$ and all pairs of
$\FamS(t)$, and $\tW$ is reached if $\We>0$.
\end{lemma}

\begin{proof}
Let $\tau^*$ be the supremum of the $t'\le T_{\mathrm{def}}$ such that the chase has not stopped before $t'$ and (J1), (J2) hold on
$[\ttd,t']$ for $\FamS$. At $\ttd$: $l=\lambda<\kappa\lambda$, and $\FamS(\ttd)$ consists of body pairs ($a_0=0<\We$ if $\We>0$),
for which (S2) gives $N<M\lambda^{1/2+\gamma}$ ($s_0<1$). So $\tau^*\ge\ttd$.

\emph{The bounds hold at $\tau^*$ \textup($<\infty$\textup).} Let $\tau^*>\ttd$. If $\tau^*\in(t_j,t_{j+1}]$, then $\Fam(t)=\Fam(\tau^*)$
for $t\in(t_j,\tau^*)$, and $\FamS(t)=\FamS(\tau^*)$ unless $\tau^*=\tW$; as $l$, $E$ are continuous in $t$
(Lemma~\ref{lem:flow}(c)), (J1), (J2) hold at $\tau^*$. If $\tau^*=\tW$, the remaining pairs are the tail pairs, and (TW), (PT) give
$l\le(1+\delta_t)\lambda<\kappa\lambda$ ($\LJ\ge1$) and $N\le s_0M\lambda^{1/2+\gamma}<M\lambda^{1/2+\gamma}$. Hence (Imp) holds at
$\tau^*$, and by (P7), (PT), (P6) its bounds are strict: $l<\kappa\lambda$, $E<\frac12M^2\lambda^{1+2\gamma}$.

\emph{Extension.} Suppose the chase does not stop at $\tau^*$. If $\tau^*$ is a leg start $t_j$, Lemma~\ref{lem:legs} (with
$t'=\tau^*$) shows that leg $j+1$ is non-final and satisfies the proviso; so in all cases the chase is defined on
$[\tau^*,\tau^*+\delta_0]$ for some $\delta_0>0$, and we may take $\delta_0$ so small that $(\tau^*,\tau^*+\delta_0]$ contains no leg
start. For $t\in(\tau^*,\tau^*+\delta_0]$, $\FamS(t)\subset\FamS(\tau^*)$ (one leg; at a leg start the family contains that of the
next leg; $\tW\notin(\tau^*,\tau^*+\delta_0]$). For a fixed pair, $l'\le(2v/\eps)l$ and $|E'|=v|\int(T-e)\cdot AT|\le(4v/\eps)l$
(Lemma~\ref{lem:excess}(i), Lemma~\ref{lem:field}(ii)). With $\Lambda:=(1+\delta_t)\LJ(U_0)$ if $\We>0$, $\Lambda:=\LJ(U_0)$ otherwise,
choose $\delta\in(0,\delta_0]$ with $e^{2v\delta/\eps}\Lambda<\kappa$ and $\delta(4v/\eps)\kappa e^{2v\delta/\eps}\le\frac12(1-\bar\theta)M^2(\eps/2)^{2\gamma}$.
Then for $t\in[\tau^*,\tau^*+\delta]$ and every pair of $\FamS(t)$ ($\lambda\ge\eps/2$):
$l_t\le e^{2v\delta/\eps}\Lambda\lambda<\kappa\lambda$ and $E(t)\le\frac12\bar\theta M^2\lambda^{1+2\gamma}+\delta(4v/\eps)\kappa
e^{2v\delta/\eps}\lambda\le\frac12M^2\lambda^{1+2\gamma}$, so \textup{(J1)}, \textup{(J2)} hold on $[\ttd,\tau^*+\delta]$, contradicting the definition of $\tau^*$.

So the chase stops at $\tau^*$, unless $\tau^*=\infty$; but then (J1) holds at all leg starts, the chase never stops and has
infinitely many legs (Lemma~\ref{lem:chase}(0)), against Lemma~\ref{lem:legs}(iv). Hence $\tau^*=\Tstop<\infty$, the proviso held at
every non-final leg start (Lemma~\ref{lem:legs}(iii)), and (J1), (J2) hold on $[\ttd,\Tstop]$. Finally, if $\We>0$, the final leg
$J+1$ has $U-a_J\le l_{t_J}(a_J,U)\le\eps$, so $a_J\ge U-\eps>\We$, and $\jW\le J$.
\end{proof}
